% Generated draft: requires mathematical and source review.
\documentclass[11pt]{article}
\usepackage[margin=25mm]{geometry}
\usepackage{fontspec}
\setmainfont{Latin Modern Roman}
\usepackage{amsmath,amssymb,mathtools,unicode-math}
\setmathfont{Latin Modern Math}
\usepackage{enumitem,needspace,xurl,hyperref}
\setlength{\parindent}{0pt}
\setlength{\parskip}{5pt}
\setlength{\emergencystretch}{3em}
\newcommand{\N}{\mathbb{N}}
\newcommand{\Z}{\mathbb{Z}}
\newcommand{\Q}{\mathbb{Q}}
\newcommand{\R}{\mathbb{R}}
\newcommand{\C}{\mathbb{C}}
\newcommand{\F}{\mathbb{F}}
\newcommand{\A}{\mathbb{A}}
\newcommand{\PP}{\mathbb{P}}
\newcommand{\E}{\mathbb{E}}
\newcommand{\eps}{\varepsilon}
\newcommand{\dd}{\,\mathrm{d}}
\begin{document}
\section*{UnsolvedMath ID: 11000077}
\textbf{Source problem code:} catalogue entry
\subsection{Definitions and mathematical statement}
Retain \(L(\operatorname{Mod}_g)\) as the source's genus-dependent quantity. Its definition is absent from the excerpt. The supplied research summary interprets it through least pseudo-Anosov dilatation. In that conventional interpretation, for closed oriented genus-\(g\) surfaces, \(g\geq2\), let \(\lambda(f)>1\) be the stretch factor of a pseudo-Anosov mapping class and set
\[
 L(\operatorname{Mod}_g)
 =\inf_{f\in\operatorname{Mod}_g\ \mathrm{pseudo\text{-}Anosov}}\log\lambda(f).
\]
This is the logarithmic convention, subject to checking the source's normalization; it is not silently identified with the unlogged dilatation or a differently normalized translation length.

\emph{Question 7.1.} Does
\[
 \lim_{g\to\infty}gL(\operatorname{Mod}_g)
\]
exist? In the ordinary finite-real-limit reading this asks whether some \(\ell\in\R\) satisfies
\[
 \forall\eps>0\ \exists g_0\ \forall g\geq g_0,
 \bigl|gL(\operatorname{Mod}_g)-\ell\bigr|<\eps.
\]
The source does not propose a value, ask only for bounds, or replace the full integer-genus sequence by a subsequence. It does not explicitly say whether an extended-real limit counts, and the definition of \(L\) and the surface convention must be recovered.

The research summary mentions Penner, a ``Penner constant,'' Tsai, and other growth bounds without precise references. It describes results using \(\log\lambda_g\sim C/g\), but no constant or theorem is supplied. A literal asymptotic equivalent of that form for one fixed \(C\) would already settle this normalized-limit question; it must not be confused with two-sided order \(1/g\) bounds or adopted as verified evidence here.
\subsection{Short English statement}
Does the genus-normalized least pseudo-Anosov logarithmic dilatation converge as genus grows? The source notation and normalization are omitted, and two-sided growth estimates alone would not establish a limit.
\subsection{Sources}
\begin{itemize}
\item Local UnsolvedMath registry, permanent ID 11000077, source code catalogue entry, Question 7.1, PDF page 57. No individual proposer is supplied.
\item Farb, \emph{Problems on Mapping Class Groups and Related Topics}. Supplied URL: \url{https://www.math.uchicago.edu/~farb/papers/mcgbook.pdf}; recorded access 2026-07-29; extraction pdf-font-structure. The edited volume contains 329 explicitly typeset headings across chapters, with repeated local numbering disambiguated by page.
\item The supplied summary interprets the question as least pseudo-Anosov dilatation, refers to Penner, a Penner constant, Tsai, and unspecified other bounds, but provides no full bibliography, exact estimates, or limit theorem.
\item Supplied status open and OPEN-TRIAGE, with the summary calling the limit subtle and not fully settled. Its potentially stronger asymptotic-equivalence wording is flagged above. Current primary-literature status and redistribution terms for the publicly accessible author/publisher copy remain NEEDS\_REVIEW.
\end{itemize}
\end{document}
